\ifdefined\gkver\else\def\gkver{student}\fi
\documentclass[\gkver]{gaokaozhenti}
\begin{document}

\chapter{2026年重庆}
\begin{enumerate}
\item
%% number: 1
%% paperName: 2026年重庆高考第 1 题 4 分
%% typeId: 1
%% type: 单选题
%% chapter: 曲线运动
%% point: 竖直平面圆周运动
%% method: 无
%% score: 4
%% degree: 850
%% duplicateId: 0
%% body:
2026 年我国自主设计和制造的摩托车在世界大赛中多次夺冠，实现了历史性突破。如图所示，某次比赛中，骑手和车（整体视为质点）以一定速度经过拱形赛道最高点时，骑手和车 \xzanswer{D}
\onepicture[w=7cm]{figs/01.png}

\fourchoices[answer=D]
{向心加速度为零}
{速度越大向心力越小}
{所受支持力大于重力}
{所受重力与支持力的合力提供向心力}

%% answer: D

%% memo:
\memoanswer{
A．质点在拱形赛道最高点，速度存在沿轨迹切线方向的分量，存在指向圆心（竖直向下）的向心加速度，向心加速度不为零，A 错误；

B．圆周运动向心力公式为 $F_{n}=m\frac{v^{2}}{r}$，在质量 $m$、轨道半径 $r$ 不变时，速度越大，向心力越大，B 错误；

CD．在最高点，向心力方向竖直向下指向圆周运动的圆心，对整体受力分析：竖直向下的重力 $mg$、竖直向上的赛道支持力 $N$，二者的合力提供竖直向下的向心力，满足 $mg-N=m\frac{v^{2}}{r}$

可得 $N=mg-m\frac{v^{2}}{r}$

因此支持力小于重力，C 错误，D 正确。

故选 D。
}
\item
%% number: 2
%% paperName: 2026年重庆高考第 2 题 4 分
%% typeId: 1
%% type: 单选题
%% chapter: 光学
%% point: 光的电磁说
%% method: 无
%% score: 4
%% degree: 780
%% duplicateId: 0
%% body:
5G 通信技术使用了厘米波段和毫米波段的电磁波，我国在此领域位于世界前列。厘米波波长在 $10^{-2} \Um$ 数量级，毫米波波长在 $10^{-3} \Um$ 数量级，则 \xzanswer{A}

\fourchoices[answer=A]
{厘米波的频率比毫米波的低}
{厘米波的光子能量比毫米波的大}
{厘米波在真空中的波速比毫米波的大}
{厘米波在真空中的波速比在介质中的小}

%% answer: A

%% memo:
\memoanswer{
电磁波在真空中传播速度均为光速 $c=3\times10^{8} \Ums$，满足关系 $c=\lambda f$，光子能量 $E=hf$（$h$ 为普朗克常量），在介质中传播速度 $v=\frac{c}{n}$（$n$ 为介质折射率，$n>1$）。

A．由 $f=\frac{c}{\lambda}$ 可知，频率与波长成反比，厘米波波长大于毫米波波长，因此厘米波频率更低，故 A 正确；

B．光子能量 $E=hf$，厘米波频率更低，因此光子能量更小，故 B 错误；

C．所有电磁波在真空中的波速均为光速 $c$，与波长无关，二者波速相等，故 C 错误；

D．由 $v=\frac{c}{n}$，$n>1$ 可得介质中波速 $v<c$，即厘米波在真空中波速比介质中大，故 D 错误。

故选 A。
}
\item
%% number: 3
%% paperName: 2026年重庆高考第 3 题 4 分
%% typeId: 1
%% type: 单选题
%% chapter: 交流电
%% point: 交流电
%% method: 无
%% score: 4
%% degree: 840
%% duplicateId: 0
%% body:
重庆夜景闻名遐迩。吊脚楼在霓虹灯映衬下呈现出梦幻般的色彩。为启动霓虹灯，将变压器的原线圈接入有效值为 $220 \UV$ 的家用交流电，副线圈接霓虹灯。当副线圈两端的电压有效值为 $15400 \UV$ 时，流过霓虹灯的电流有效值为 $10 \UmA$，若将该变压器视为理想变压器，则 \xzanswer{B}

\fourchoices[answer=B]
{原、副线圈匝数比为 $70:1$}
{原线圈中电流有效值为 $0.7 \UA$}
{家用交流电的电压最大值为 $110\sqrt{2} \UV$}
{副线圈的输出功率为 $2.2 \UW$}

%% answer: B

%% memo:
\memoanswer{
已知理想变压器原线圈电压有效值 $U_{1}=220 \UV$，副线圈电压有效值 $U_{2}=15400 \UV$，副线圈电流 $I_{2}=10 \UmA=0.01 \UA$，结合理想变压器规律逐一分析选项：

A．理想变压器电压与匝数成正比，即 $\frac{n_{1}}{n_{2}}=\frac{U_{1}}{U_{2}}=\frac{220}{15400}=\frac{1}{70}$，原、副线圈匝数比为 $1:70$，故 A 错误；

B．理想变压器输入功率等于输出功率，副线圈输出功率 $P_{2}=U_{2}I_{2}=15400\times0.01 \UW=154 \UW$，原线圈电流 $I_{1}=\frac{P_{1}}{U_{1}}=\frac{P_{2}}{U_{1}}=\frac{154}{220} \UA=0.7 \UA$，故 B 正确；

C．家用正弦交流电的最大值为有效值的 $\sqrt{2}$ 倍，有效值为 $220 \UV$，因此最大值为 $220\sqrt{2} \UV$，故 C 错误；

D．由上述计算可知副线圈输出功率为 $154 \UW$，故 D 错误。

故选 B。
}
\item
%% number: 4
%% paperName: 2026年重庆高考第 4 题 4 分
%% typeId: 1
%% type: 单选题
%% chapter: 曲线运动
%% point: 圆周运动的应用
%% method: 无
%% score: 4
%% degree: 660
%% duplicateId: 0
%% body:
马年春晚《喜雨》舞蹈用珠帘勾勒出“春雨万物生”的唯美景象。小明用如图所示的模型分析珠帘上珠子的运动，他用柔软轻质细线连接两个相同的小球 $b_{1}$、$b_{2}$，当两小球以相同角速度绕竖直轴 $OO'$ 做匀速圆周运动时，细线①、②与竖直方向夹角分别为 $\theta_{1}$ 和 $\theta_{2}$（$\theta_{1}<\theta_{2}$），若细线①、②上弹力大小分别为 $F_{1}$ 和 $F_{2}$，$b_{1}$、$b_{2}$ 的线速度大小分别为 $v_{1}$ 和 $v_{2}$，向心加速度大小分别为 $a_{1}$ 和 $a_{2}$，忽略空气阻力，则 \xzanswer{D}
\onepicture[h=5.0cm]{figs/04.png}

\fourchoices[answer=D]
{$v_{1}>v_{2}$}
{$2a_{2}>a_{1}>a_{2}$}
{$F_{1}>2F_{2}$}
{$2F_{2}>F_{1}>F_{2}$}

%% answer: D

%% memo:
\memoanswer{
A．匀速圆周运动线速度满足 $v=\omega r$，$b_{2}$ 的圆周运动半径 $r_{2}$ 是 $b_{1}$ 的半径 $r_{1}$ 加上线②的水平投影，因此 $r_{2}>r_{1}$，可得 $v_{2}=\omega r_{2}>v_{1}=\omega r_{1}$，A 错误；

B．向心加速度满足 $a=\omega^{2}r$，因 $r_{2}>r_{1}$，故 $a_{2}=\omega^{2}r_{2}>a_{1}=\omega^{2}r_{1}$，B 错误；

CD．对下方小球，竖直方向合力为 0，$F_{2}\cos\theta_{2}=mg$，得 $F_{2}=\frac{mg}{\cos\theta_{2}}$

对两球整体受力分析，竖直方向合力为 0，$F_{1}\cos\theta_{1}=2mg$，得 $F_{1}=\frac{2mg}{\cos\theta_{1}}$

已知 $\theta_{1}<\theta_{2}$，则 $\cos\theta_{1}>\cos\theta_{2}$，因此 $\frac{1}{\cos\theta_{1}}<\frac{1}{\cos\theta_{2}}$

代入得 $F_{1}=\frac{2mg}{\cos\theta_{1}}<\frac{2mg}{\cos\theta_{2}}=2F_{2}$

由两球水平方向向心力关系 $F_{1}\sin\theta_{1}-F_{2}\sin\theta_{2}=ma_{1}$，$F_{2}\sin\theta_{2}=ma_{2}$

由此可知 $F_{1}\sin\theta_{1}>F_{2}\sin\theta_{2}$

同时 $\sin\theta_{1}<\sin\theta_{2}$，故 $F_{1}>F_{2}$

最终得 $2F_{2}>F_{1}>F_{2}$，C 错误，D 正确。

故选 D。
}
\item
%% number: 5
%% paperName: 2026年重庆高考第 5 题 4 分
%% typeId: 1
%% type: 单选题
%% chapter: 光学
%% point: 光的折射
%% method: 无
%% score: 4
%% degree: 550
%% duplicateId: 0
%% body:
某种特殊眼镜通过三棱镜改变光路给人们视物带来方便。如图所示，三棱镜的 $MN$ 面为镜面，光从空气经 $MK$ 面射入棱镜，最终从 $NK$ 面射出，进入眼睛。若 $\angle M=30^{\circ}$，某单色光从 $MK$ 面 $P$ 点与 $MK$ 成 $45^{\circ}$ 角射入后，垂直射到镜面 $MN$ 上，则 \xzanswer{C}
\onepicture[h=4.6cm]{figs/05.png}

\fourchoices[answer=C]
{三棱镜对该光的折射率为 $1.5$}
{三棱镜对该光的折射率为 $\sqrt{3}$}
{该光从三棱镜射入空气发生全反射的临界角为 $45^{\circ}$}
{该光从三棱镜射入空气发生全反射的临界角为 $60^{\circ}$}

%% answer: C

%% memo:
\memoanswer{
AB．入射光线与 $MK$ 界面成 $45^{\circ}$ 角射入，入射角是入射光线与界面法线的夹角，因此空气中的入射角：$i=90^{\circ}-45^{\circ}=45^{\circ}$

已知光进入棱镜后垂直射到镜面 $MN$ 上，可得棱镜内的折射角：$r=90^{\circ}-60^{\circ}=30^{\circ}$

根据折射定律 $n=\frac{\sin i}{\sin r}$

代入数值：$n=\frac{\sin45^{\circ}}{\sin30^{\circ}}=\sqrt{2}$，故 AB 错误；

CD．全反射临界角 $C$ 满足 $\sin C=\frac{1}{n}$，代入得：$\sin C=\frac{1}{\sqrt{2}}\Rightarrow C=45^{\circ}$，故 C 正确，D 错误。

故选 C。
}
\item
%% number: 6
%% paperName: 2026年重庆高考第 6 题 4 分
%% typeId: 1
%% type: 单选题
%% chapter: 光学
%% point: 波粒二象性
%% method: 无
%% score: 4
%% degree: 630
%% duplicateId: 0
%% body:
科学研究中，通过“角分辨光电子能谱仪”测量光电效应中光电子的动能和出射角，可了解样品的相关特性。若一束频率为 $\nu$ 的紫外光照射在样品表面产生光电子，某具有最大初动能 $E_{k}$ 的光电子，其速度方向与样品表面法线的夹角为 $\theta$，光电子质量为 $m_{c}$，普朗克常量为 $h$，则 \xzanswer{B}

\fourchoices[answer=B]
{样品的逸出功为 $h\nu-E_{k}\cos\theta$}
{该光电子的德布罗意波长为 $\frac{h}{\sqrt{2m_{c}E_{k}}}$}
{该光电子的速度沿样品表面法线方向的分量大小为 $\sqrt{\frac{2E_{k}}{m_{c}}}\sin\theta$}
{当入射光的频率变为 $2\nu$ 时，光电子的最大初动能变为 $2E_{k}$}

%% answer: B

%% memo:
\memoanswer{
A．根据爱因斯坦光电效应方程 $E_{k}=h\nu-W_{0}$，光电子的最大初动能 $E_{k}$ 仅与入射光频率和逸出功有关，和出射方向无关，因此逸出功 $W_{0}=h\nu-E_{k}$，故 A 错误；

B．德布罗意波长满足 $\lambda=\frac{h}{p}$，光电子的动能与动量关系为 $E_{k}=\frac{p^{2}}{2m_{c}}$

联立得 $\lambda=\frac{h}{\sqrt{2m_{c}E_{k}}}$，故 B 正确；

C．光电子速度方向与样品表面法线夹角为 $\theta$，因此法线方向速度分量为 $v\cos\theta$，结合 $E_{k}=\frac{1}{2}m_{c}v^{2}$ 得 $v=\sqrt{\frac{2E_{k}}{m_{c}}}$

故法线分量大小应为 $\sqrt{\frac{2E_{k}}{m_{c}}}\cos\theta$，而选项给出的是 $\sin\theta$，故 C 错误；

D．当入射光频率变为 $2\nu$ 时，新的最大初动能 $E_{k}'=2h\nu-W_{0}$

代入 $W_{0}=h\nu-E_{k}$ 得 $E_{k}'=h\nu+E_{k}\neq2E_{k}$，故 D 错误。

故选 B。
}
\item
%% number: 7
%% paperName: 2026年重庆高考第 7 题 4 分
%% typeId: 1
%% type: 单选题
%% chapter: 磁场
%% point: 带电粒子在磁场中的运动
%% method: 无
%% score: 4
%% degree: 310
%% duplicateId: 0
%% body:
某同学设计了一种测量粒子速率的装置，其原理示意图如图所示。空间充满匀强磁场，粒子源位于 $K$ 点，可发射比荷相同、速率分别为 $v_{1}$ 和 $v_{2}$ 的带正电粒子（$v_{1}<v_{2}$）。粒子源右侧有一过 $C$ 点的竖直挡板。若磁场方向垂直于纸面向外，水平向右发射的粒子打在 $M$、$H$ 两点；若磁场方向垂直于纸面向里，水平向左发射的粒子打在 $H$、$N$ 两点，$CM:MN=1:8$，$KC\perp CN$，不考虑粒子重力及粒子间相互作用，则 $v_{1}:v_{2}$ 为 \xzanswer{A}
\onepicture[h=5.2cm]{\includesvg{figs/07}}

\fourchoices[answer=A]
{$3:5$}
{$2:3$}
{$3:4$}
{$4:5$}

%% answer: A

%% memo:
\memoanswer{
带电粒子在匀强磁场中做匀速圆周运动，洛伦兹力提供向心力，由 $qvB=m\frac{v^{2}}{r}$

可得轨道半径公式：$r=\frac{mv}{qB}$

由于粒子比荷 $\frac{q}{m}$ 相同、磁场 $B$ 固定，因此轨道半径与速率成正比，即 $r\propto v$

根据题述信息，以 $K$ 为上顶点作出两圆，此两圆与直线 $MN$ 仅有 3 个交点，如图所示。

\onepicture[h=4.2cm]{figs/07_an.jpg}

磁场垂直纸面向外、粒子水平向右发射时：带正电粒子受洛伦兹力向下，圆心在 $K$ 点正下方。速率小的 $v_{1}$ 对应轨道半径 $r_{1}$，打在 $H$ 点；速率大的 $v_{2}$ 对应轨道半径 $r_{2}$，打在 $M$ 点；

磁场垂直纸面向里、粒子水平向左发射时：带正电粒子受洛伦兹力向下，圆心在 $K$ 点正下方。速率小的 $v_{1}$ 打在 $H$ 点，速率大的 $v_{2}$ 打在 $N$ 点；

由几何关系可得：$r_{1}=CK=CH$

由几何关系可得：

\[ r_{2}^{2}=\left(\frac{MN}{2}\right)^{2}+\left(CK\right)^{2},\quad 2CM=2r_{2}-MN \]

结合题设条件 $CM:MN=1:8$

联立可推导出 $v_{1}:v_{2}=r_{1}:r_{2}=3:5$

故选 A。
}
\item
%% number: 8
%% paperName: 2026年重庆高考第 8 题 5 分
%% typeId: 2
%% type: 多选题
%% chapter: 曲线运动
%% point: 人造地球卫星
%% method: 无
%% score: 5
%% degree: 550
%% duplicateId: 0
%% body:
中国空间站已成为我国规模最大、功能最全的综合性近地空间研究平台。若空间站绕地球的运动视为匀速圆周运动，轨道半径为 $r$，地球半径为 $R$，地球表面的重力加速度大小为 $g$，引力常量为 $G$，忽略地球自转的影响，则 \xzanswer{CD}

\fourchoices[answer=CD]
{地球的质量为 $\frac{gr^{2}}{G}$}
{空间站绕地球运行的向心加速度大小大于 $g$}
{空间站绕地球运行的周期为 $2\pi\sqrt{\frac{r^{3}}{gR^{2}}}$}
{空间站绕地球运行的线速度大小为 $\sqrt{\frac{gR^{2}}{r}}$}

%% answer: CD

%% memo:
\memoanswer{
A．忽略地球自转的影响，地球表面的物体的重力等于万有引力有 $\frac{GMm}{R^{2}}=mg$ 解得地球的质量为 $M=\frac{gR^{2}}{G}$，故 A 错误；

B．由万有引力提供向心力有 $\frac{GMm'}{r^{2}}=m'a$ 解得 $a=\frac{GM}{r^{2}}$

代入地球质量可得 $a=\frac{R^{2}}{r^{2}}g$

由于 $R<r$，则有 $a<g$，故 B 错误；

C．由万有引力提供向心力有 $\frac{GMm'}{r^{2}}=m'\frac{4\pi^{2}}{T^{2}}r$ 解得 $T=\sqrt{\frac{4\pi^{2}r^{3}}{GM}}$

代入地球质量可得 $T=2\pi\sqrt{\frac{r^{3}}{gR^{2}}}$，故 C 正确；

D．由万有引力提供向心力有 $\frac{GMm'}{r^{2}}=m'\frac{v^{2}}{r}$ 解得 $v=\sqrt{\frac{GM}{r}}$

代入地球质量可得 $v=\sqrt{\frac{gR^{2}}{r}}$，故 D 正确；

故选 CD。
}
\item
%% number: 9
%% paperName: 2026年重庆高考第 9 题 5 分
%% typeId: 2
%% type: 多选题
%% chapter: 电路
%% point: 闭合电路欧姆定律
%% method: 无
%% score: 5
%% degree: 630
%% duplicateId: 0
%% body:
图~\ref{2026重庆09a} 为某同学设计的可调亮度台灯的电路图。电源 $E$ 的电动势为 $6 \UV$，内阻为 $1 \UO$，灯泡 L 和三个电阻的阻值均为 $R$。$S_{1}$ 闭合，当 $S_{2}$ 处于图~\ref{2026重庆09a} 所示断开状态时，测得 L 两端电压为 $1.80 \UV$，再将 $S_{2}$ 接 $a$ 或 $b$，测量 L 两端电压。若 L 阻值视为不变，则 \xzanswer{AC}
\twopicture[numA=1, numB=2, labelA=2026重庆09a, labelB=2026重庆09b, h=3.6cm]
{figs/09a.png}
{figs/09b.png}

\fourchoices[answer=AC]
{$R$ 为 $3 \UO$}
{图~\ref{2026重庆09b} 示数为 $S_{2}$ 接 $a$ 时 L 两端电压}
{图~\ref{2026重庆09b} 示数为 $S_{2}$ 接 $b$ 时 L 两端电压}
{$S_{2}$ 接 $b$ 时外电路的电功率约为 $3.74 \UW$}

%% answer: AC

%% memo:
\memoanswer{
A．当 $S_{1}$ 闭合、$S_{2}$ 断开时，电路总结构为：灯泡 L 与右侧两个串联的电阻 $R$ 串联，再和电源内阻 $r=1 \UO$ 串联，总电阻为：$R_{\text{总}}=R+R+R+r=3R+r$

已知灯泡两端电压 $U_{L}=1.8 \UV$，由欧姆定律得总电流 $I=\frac{U_{L}}{R}=\frac{1.8}{R}$

结合闭合电路欧姆定律 $E=IR_{\text{总}}$

代入 $E=6 \UV$ 得 $R=3 \UO$，故 A 正确；

BC．$S_{2}$ 接 $a$ 时：并联部分为两个阻值为 $R$ 的电阻并联，并联总电阻 $R_{\text{并}a}=\frac{R\cdot R}{R+R}=1.5 \UO$

电路总电阻 $R_{\text{总}a}=R_{L}+R_{\text{并}a}+R+r=8.5 \UO$

总电流 $I_{a}=\frac{E}{R_{\text{总}a}}=\frac{6}{8.5} \UA$

灯泡两端电压：$U_{La}=I_{a}\cdot R_{L}=\frac{6}{8.5}\times3 \UV\approx2.12 \UV$

$S_{2}$ 接 $b$ 时：仅有灯泡和一个电阻串联接入电源，电路总电阻 $R_{\text{总}b}=R_{L}+R+r=7 \UO$

总电流 $I_{b}=\frac{E}{R_{\text{总}b}}=\frac{6}{7} \UA$

灯泡两端电压：$U_{Lb}=I_{b}\cdot R_{L}=\frac{6}{7}\times3 \UV\approx2.57 \UV$

观察图 2 电压表示数约为 $2.57 \UV$，显然接 $b$，故 B 错误，C 正确；

D．$S_{2}$ 接 $b$ 时，外电路电功率：$P_{\text{外}}=I_{b}^{2}(R_{L}+R)\approx4.41 \UW$，故 D 错误。

故选 AC。
}
\item
%% number: 10
%% paperName: 2026年重庆高考第 10 题 5 分
%% typeId: 2
%% type: 多选题
%% chapter: 电场
%% point: 带电粒子的偏转
%% method: 无
%% score: 5
%% degree: 420
%% duplicateId: 0
%% body:
如图所示，$xOy$ 平面内充满沿 $y$ 轴正方向的匀强电场，原点 $O$ 处的粒子源以不同的初速度沿 $x$ 轴发射多个比荷为 $k$ 的带正电粒子。将粒子在电场中运动速率相同的位置连接起来，得到一条平滑闭合曲线。图中曲线 $a$、$b$ 对应的速率分别为 $v$、$\sqrt{2}v$，直线 $y=x$ 与 $a$、$b$ 交于 $M$、$N$ 两点。若 $a$ 与 $y$ 轴交于 $(0,L)$ 点，不考虑粒子重力及粒子间相互作用，则 \xzanswer{AB}
\onepicture[h=4.6cm]{figs/10.png}

\fourchoices[answer=AB]
{电场强度的大小为 $\frac{v^{2}}{2kL}$}
{$M$、$N$ 两点间的电势差为 $\frac{2v^{2}}{5k}$}
{经过 $a$ 最右端的粒子初速度大小为 $\frac{1}{2}v$}
{粒子从 $O$ 到 $M$ 所用时间与从 $O$ 到 $N$ 所用时间之比为 $1:2$}

%% answer: AB

%% memo:
\memoanswer{
A．带正电粒子从 $O$ 点沿 $x$ 轴发射，做类平抛运动：$x$ 方向匀速 $x=v_{0}t$，$y$ 方向匀加速 $y=\frac{1}{2}at^{2}$，其中加速度 $a=\frac{qE}{m}=kE$

由动能定理，对末速率为 $v'$ 的粒子有：$qEy=\frac{1}{2}mv'^{2}-\frac{1}{2}mv_{0}^{2}\Rightarrow v'^{2}=v_{0}^{2}+2kEy$

曲线 $a$ 对应末速率 $v$，过 $(0,L)$ 点，该点粒子初速度 $v_{0}=0$，代入动能定理得 $\frac{1}{2}mv^{2}=qEL$

结合比荷 $k=\frac{q}{m}$，解得 $E=\frac{v^{2}}{2kL}$，故 A 正确；

B．$M$ 点在 $y=x$ 上且属于曲线 $a$，联立 $x=y=v_{0}t$ 和 $y=\frac{1}{2}at^{2}$ 得 $t=\frac{2v_{0}}{a}$

代入速率关系 $v^{2}=v_{0}^{2}+(at)^{2}$ 得 $v^{2}=5v_{0}^{2}$

因此 $y_{M}=\frac{2v^{2}}{5a}$

同理 $N$ 点属于曲线 $b$（末速率 $\sqrt{2}v$），得 $y_{N}=\frac{4v^{2}}{5a}$

匀强电场沿 $y$ 方向，电势差 $U_{MN}=E(y_{N}-y_{M})$

得 $U_{MN}=\frac{2v^{2}}{5k}$，故 B 正确；

C．曲线 $a$ 对应 $v^{2}=v_{0}^{2}+2kEy$

经过 $a$ 最右端时 $y=\frac{L}{2}$

整理可得 $v^{2}=v_{0}^{2}+kEL=v_{0}^{2}+\frac{v^{2}}{2}$

解得 $v_{0}=\frac{\sqrt{2}}{2}v$，故 C 错误；

D．由 $y=\frac{1}{2}kEt^{2}$ 得 $t=\sqrt{\frac{2y}{kE}}$

代入 $y_{M}=\frac{2v^{2}}{5a}$、$y_{N}=\frac{4v^{2}}{5a}$，可得 $\frac{t_{M}}{t_{N}}=\sqrt{\frac{y_{M}}{y_{N}}}=\sqrt{\frac{1}{2}}$，故 D 错误。

故选 AB。
}
\item
%% number: 11
%% paperName: 2026年重庆高考第 11 题 7 分
%% typeId: 4
%% type: 实验题
%% chapter: 力物体的平衡
%% point: 胡克定律
%% method: 无
%% score: 7
%% degree: 480
%% duplicateId: 0
%% body:
气弹簧常用于生产生活中许多减震系统，通常由密闭有气体的气缸、活塞和活塞杆组成。气弹簧中活塞杆的一端位于气缸内、并固定有活塞，另一端位于气缸外。某兴趣小组通过实验，测试了气弹簧的特性。根据下列信息完成以下小题。

他们研究活塞杆伸出的长度与缸内气体对活塞压力之间的关系。实验装置如图~\ref{2026重庆11a} 所示，气弹簧竖直放置，上端固定于水平横梁，下端连接测力装置。
\twopicture[numA=1, numB=2, labelA=2026重庆11a, labelB=2026重庆11b, h=4.4cm]
{figs/11a.png}
{figs/11b.png}
\begin{enumerate}
	\item
	根据实验设计，要求测量活塞杆伸出长度 $x$ 的仪器最小分度值为 $1 \Umm$，则最合适的测量仪器为 \tkanswer{毫米刻度尺}（填“毫米刻度尺”或“螺旋测微器”）。
	\item
	根据实验数据得到缸内气体对活塞的压力 $F$ 随 $x$ 的变化关系如图~\ref{2026重庆11b} 所示（已消除大气压力、摩擦力、活塞和活塞杆质量的影响）。若将气弹簧等效为弹簧，则由图可知，其劲度系数约为 \tkanswer{$367$} $\UNm$（保留 3 位有效数字）。
	\item
	由图~\ref{2026重庆11b} 可知，$x$ 从 $40.0 \Umm$ 变化到 $70.0 \Umm$ 过程中，缸内气体对外做功约为 \tkanswer{$3.86$} $\UJ$（保留 3 位有效数字）。
\end{enumerate}

%% answer:
\jdanswer{
\begin{enumerate}
	\item
	毫米刻度尺

	\item
	$367 \UNm$

	\item
	$3.86 \UJ$
\end{enumerate}
}
%% memo:
\memoanswer{
（1）毫米刻度尺的最小分度值为 $1 \Umm$，符合测量要求，螺旋测微器最小分度值为 $0.01 \Umm$，精度远高于需求，因此选毫米刻度尺。

（2）等效劲度系数满足 $k=\left|\frac{\Delta F}{\Delta x}\right|$

从 $F-x$ 图取直线上两点，例如 $x_{1}=40.0 \Umm$，$F_{1}=134.0 \UN$；$x_{2}=70.0 \Umm$，$F_{2}=123.0 \UN$

代入得 $k=\frac{|123.0-134.0|}{(70.0-40.0)\times10^{-3}} \UNm\approx367 \UNm$

（3）$F-x$ 图线与 $x$ 轴围成的面积对应气体对外做功，该过程 $F$ 近似线性变化，用平均力计算：$W=\overline{F}\cdot\Delta x=\frac{F_{1}+F_{2}}{2}\cdot\Delta x=\frac{134.0+123.0}{2}\times30\times10^{-3} \UJ\approx3.86 \UJ$
}
\item
%% number: 12
%% paperName: 2026年重庆高考第 12 题 9 分
%% typeId: 4
%% type: 实验题
%% chapter: 气体性质
%% point: 玻意耳定律
%% method: 无
%% score: 9
%% degree: 370
%% duplicateId: 0
%% body:
气弹簧常用于生产生活中许多减震系统，通常由密闭有气体的气缸、活塞和活塞杆组成。气弹簧中活塞杆的一端位于气缸内、并固定有活塞，另一端位于气缸外。某兴趣小组通过实验，测试了气弹簧的特性。根据下列信息完成以下小题。

他们研究了气弹簧伸缩过程中气缸内密闭气体的状态变化。
\twopicture[numA=1, numB=2, labelA=2026重庆12a, labelB=2026重庆12b, h=3.8cm]
{figs/12a.png}
{figs/12b.png}
\begin{enumerate}
	\item
	用 50 分度游标卡尺测量活塞杆的直径示数如图~\ref{2026重庆12a} 所示，则活塞杆直径 $d$ 的读数为 \tkanswer{$6.04$} $\Umm$。
	\item
	气缸的结构示意图如图~\ref{2026重庆12b} 所示。气缸被活塞分成 1 和 2 两个互相连通的气室，达到平衡时，两个气室内气体在活塞上、下表面的压力差就是缸内气体对活塞的压力 $F$，缸内气体的压强 $p=$ \tkanswer{$\frac{4F}{\pi d^{2}}$}（用 $F$ 和 $d$ 表示）。
	\item
	缸内气体的最小体积为 $V_{0}$，忽略活塞厚度，活塞杆的总长度与气缸的长度视为相同，活塞杆伸出长度为 $x$ 时缸内气体的体积 $V=$ \tkanswer{$V_{0}+\frac{\pi d^{2}x}{4}$}（用 $V_{0}$、$d$ 和 $x$ 表示）。
	\item
	由实验数据得到 $F$ 随 $\frac{1}{V}$ 变化的拟合直线如图~\ref{2026重庆12c} 所示。该小组在拟合直线上读取两点坐标（$10.40\times10^{4} \Umnc$，$134.3 \UN$）和（$9.40\times10^{4} \Umnc$，$121.3 \UN$），请写出拟合直线方程，由此判断缸内气体的状态变化是否符合理想气体的等温变化规律并给出理由。
\end{enumerate}
\onepicture[num=3, label=2026重庆12c, w=9.5cm]{\includesvg{figs/12c}}

%% answer:
\jdanswer{
\begin{enumerate}
	\item
	$6.04 \Umm$

	\item
	$\frac{4F}{\pi d^{2}}$

	\item
	$V_{0}+\frac{\pi d^{2}x}{4}$

	\item
	$F=1.30\times10^{-3}\frac{1}{V}-0.9 \UN$；不符合，理由见详解。
\end{enumerate}
}
%% memo:
\memoanswer{
（1）50 分度游标卡尺精度为 $0.02 \Umm$，主尺读数为 $6 \Umm$，游标尺第 2 条刻度线与主尺对齐，游标尺读数 $2\times0.02 \Umm=0.04 \Umm$

总读数为 $6 \Umm+0.04 \Umm=6.04 \Umm$

（2）压力差 $F$ 是气体压强作用在活塞杆横截面上的合力，活塞杆横截面积 $S=\frac{\pi d^{2}}{4}$。由 $F=pS$ 变形得 $p=\frac{F}{S}=\frac{4F}{\pi d^{2}}$

（3）活塞杆伸出长度为 $x$ 时，原本位于气缸内长度为 $x$ 的活塞杆移出气缸，缸内气体新增的体积等于这段活塞杆的体积，即 $\Delta V=Sx=\frac{\pi d^{2}x}{4}$

因此总体积为最小体积加上新增体积：$V=V_{0}+\frac{\pi d^{2}x}{4}$

（4）拟合直线斜率：$k'=\frac{\Delta F}{\Delta(1/V)}=\frac{134.3-121.3}{10.40\times10^{4}-9.40\times10^{4}} \UN\cdot\Umc=1.3\times10^{-3} \UN\cdot\Umc$

代入一点求截距：$121.3 \UN=1.30\times10^{-3}\times9.40\times10^{4} \UN+b$

计算得截距 $b=-0.9 \UN$

因此实际拟合直线方程为：$F=1.30\times10^{-3}\frac{1}{V}-0.9 \UN$（单位：$F$ 为 $\UN$，$\frac{1}{V}$ 为 $\Umnc$）

判断：不符合理想气体等温变化规律。

理由：理想气体等温变化满足玻意耳定律 $pV=C$，代入 $p=\frac{4F}{\pi d^{2}}$

可得 $F=\frac{\pi d^{2}C}{4}\cdot\frac{1}{V}$

即 $F$ 与 $\frac{1}{V}$ 呈正比例过原点的线性关系，本次拟合的直线不过原点，不符合等温变化的特征。
}
\item
%% number: 13
%% paperName: 2026年重庆高考第 13 题 10 分
%% typeId: 6
%% type: 计算题
%% chapter: 运动和力
%% point: 牛顿定律的应用
%% method: 无
%% score: 10
%% degree: 680
%% duplicateId: 0
%% body:
《天工开物》记载了一种传统木工工具——刨子。如图所示，刨子主要由刨身和刨刀构成。为体验使用刨子，小明先将刨刀取下，用大小为 $F$、方向斜向下并与水平面夹角为 $\theta$ 的力作用于刨身，使其在木材水平表面做匀速直线运动，已知重力加速度大小为 $g$。
\begin{enumerate}
	\item
	求刨身做匀速直线运动时受到的滑动摩擦力大小。
	\item
	若刨身与木材间的动摩擦因数为 $\mu$，求刨身的质量。
	\item
	将刨刀装入后，刨子质量为 $M$，改变推力，使刨子在木材水平表面由静止开始做匀加速直线运动。若加速度大小为 $a$，运动距离为 $L$，求该过程中刨子所受合力的冲量大小。
\end{enumerate}
\onepicture[w=7cm, align=right]{figs/13.png}

%% answer: （1）$F\cos\theta$；（2）$\frac{F(\cos\theta-\mu\sin\theta)}{\mu g}$；（3）$M\sqrt{2aL}$
\jdanswer{
\begin{enumerate}
	\item
	$F\cos\theta$

	\item
	$\frac{F(\cos\theta-\mu\sin\theta)}{\mu g}$

	\item
	$M\sqrt{2aL}$
\end{enumerate}
}
%% memo:
\memoanswer{
（1）刨身匀速直线运动，水平方向受力平衡。将斜向下的推力 $F$ 沿水平、竖直方向分解，水平分力为 $F\cos\theta$，滑动摩擦力与水平分力大小相等，因此滑动摩擦力 $f=F\cos\theta$

（2）竖直方向受力平衡，木材对刨身的支持力 $N$ 满足 $N=mg+F\sin\theta$

滑动摩擦力满足 $f=\mu N$

其中 $f=F\cos\theta$

代入得 $F\cos\theta=\mu(mg+F\sin\theta)$

整理后解得刨身质量 $m=\frac{F(\cos\theta-\mu\sin\theta)}{\mu g}$

（3）刨子由静止开始做匀加速直线运动，由运动学公式 $v^{2}=2aL$

可得末速度 $v=\sqrt{2aL}$

根据动量定理，合力的冲量等于物体动量的变化量，初动量为 0，因此合力冲量 $I=\Delta p=Mv-0=M\sqrt{2aL}$
}
\item
%% number: 14
%% paperName: 2026年重庆高考第 14 题 13 分
%% typeId: 6
%% type: 计算题
%% chapter: 直线运动
%% point: 多段匀变速直线运动
%% method: 无
%% score: 13
%% degree: 640
%% duplicateId: 0
%% body:
智能交通控制中，为提升通行效率，某些路段设置了“绿波带”，能使以一定速度行驶的车辆连续绿灯通行。小明选取某平直道路相邻三个路口 $A$、$B$、$C$ 的信号灯进行模拟研究（如图），路口间距 $AB=BC=600 \Um$。三个路口绿灯、红灯持续时间均为 $30 \Us$。$t=0$ 时，$A$ 的绿灯开启；$t=30 \Us$ 时，$B$ 的绿灯开启；$t=90 \Us$ 时，$C$ 的绿灯开启。$t=0$ 时，车由静止从 $A$ 开始做匀加速直线运动直到速度大小为 $15 \Ums$，之后一直做匀速直线运动并依次通过 $B$、$C$ 路口。车视为质点，不计通过路口的时间，不考虑黄灯时间。
\begin{enumerate}
	\item
	求车做匀加速直线运动的平均速度大小。
	\item
	若 $t=50 \Us$ 时，车刚好通过 $B$ 路口，求车做匀加速直线运动的位移大小。
	\item
	若车在 $90 \Us\leq t\leq100 \Us$ 内通过 $C$ 路口，求车做匀加速直线运动的加速度大小范围。
\end{enumerate}
\onepicture[w=10cm, align=right]{\includesvg{figs/14}}

%% answer: （1）$7.5\Ums$；（2）$150\Um$；（3）$0.375\Umsq\leq a\leq0.75\Umsq$
\jdanswer{
\begin{enumerate}
	\item
	$7.5 \Ums$

	\item
	$150 \Um$

	\item
	$0.375 \Umsq\leq a\leq0.75 \Umsq$
\end{enumerate}
}
%% memo:
\memoanswer{
（1）对于匀变速直线运动，某段过程的平均速度满足 $\overline{v}=\frac{v_{0}+v_{t}}{2}$

汽车由静止开始匀加速，初速度 $v_{0}=0$，末速度 $v=15 \Ums$，因此平均速度：$\overline{v}=\frac{0+15}{2} \Ums=7.5 \Ums$

（2）汽车从 $A$ 到 $B$ 总位移 $x_{AB}=600 \Um$，总用时 $t_{1}=50 \Us$。设匀加速时间为 $t_{0}$，则匀速运动的时间为 $t_{1}-t_{0}$

总位移可表示为匀加速位移加匀速位移：$x_{AB}=\overline{v}\cdot t_{0}+v\cdot(t_{1}-t_{0})$

解得 $t_{0}=20 \Us$

因此匀加速位移：$x_{\text{加}}=\overline{v}\cdot t_{0}=150 \Um$

（3）汽车从 $A$ 到 $C$ 总位移 $x_{AC}=1200 \Um$，要求总运动时间 $90 \Us\leq t\leq100 \Us$。设匀加速加速度为 $a$，则匀加速到 $15 \Ums$ 的时间 $t_{0}=\frac{v}{a}=\frac{15}{a}$

匀加速位移 $x_{\text{加}}=\frac{1}{2}at_{0}^{2}=\frac{225}{2a}$

总位移满足：$x_{AC}=x_{\text{加}}+v\cdot(t-t_{0})$

代入数值整理得：

\[ 1200=\frac{225}{2a}+15\left(t-\frac{15}{a}\right) \]

化简后得到总时间与加速度的关系：$t=80+\frac{7.5}{a}$

当 $t_{\max}=100 \Us$ 时，对应加速度最小值：$a_{\min}=0.375 \Umsq$

当 $t_{\min}=90 \Us$ 时，对应加速度最大值：$a_{\max}=0.75 \Umsq$

校验通过 $B$ 灯时，是否闯红灯：

加速到最大速度时，最大位移 $\frac{v^{2}}{2a_{\min}}=300 \Um$

说明到 $B$ 之前已达到最大速度；

车通过 $B$ 的时刻 $t_{B}$ 满足 $AB=600 \Um$，同理可得 $t_{B}$ 与加速度的关系：

\[ \frac{225}{2a}+15\left(t_{B}-\frac{15}{a}\right)=600\Rightarrow t_{B}=40+\frac{7.5}{a} \]

$B$ 的绿灯窗口为 $t\in[30,60] \Us$

结合得到的 $a\in[0.375,0.75] \Umsq$，代入 $t_{B}$ 的表达式可得：

当 $a=0.375 \Umsq$ 时，$t_{B}=60 \Us$；当 $a=0.75 \Umsq$ 时，$t_{B}=50 \Us$。

即满足 $B$ 路口绿灯通行的要求，不存在闯红灯的情况。

因此加速度范围为 $0.375 \Umsq\leq a\leq0.75 \Umsq$
}

\gknewpage
\item
%% number: 15
%% paperName: 2026年重庆高考第 15 题 16 分
%% typeId: 6
%% type: 计算题
%% chapter: 电磁感应
%% point: 电磁感应中的变加速
%% method: 分情况讨论
%% score: 16
%% degree: 190
%% duplicateId: 0
%% body:
如图甲所示，在水平面内，相距为 $L=1 \Um$ 的固定平行光滑金属导轨与滑动变阻器 $R$ 相连，滑动变阻器的最大阻值为 $3 \UO$，导轨电阻不计。空间内存在磁感应强度大小为 $B=2 \UT$、方向竖直向下的匀强磁场。粗细均匀的导体棒在外力 $F$ 作用下以初速度 $v_{0}$ 向右运动，导体棒长为 $L$，质量 $m=2 \Ukg$ 电阻不计。运动过程中，外力 $F$ 随金属棒的动能 $E_{k}$ 变化图像如图乙所示。金属棒和导轨接触良好且始终和导轨垂直，忽略回路电流产生的磁场。
\begin{enumerate}
	\item
	滑动变阻器阻值 $R=1 \UO$，求金属棒速度为 $0.5 \Ums$ 时所受的安培力大小；
	\item
	调节滑动变阻器阻值 $R=1 \UO$，初速度 $v_{0}=\frac{5}{6} \Ums$，求金属棒的稳定速度；当金属棒以不同初速度 $v_{0}$ 运动时，求在达到稳定速度的过程中合外力做功；
	\item
	调节滑动变阻器阻值，使得金属棒存在两个稳定速度，求滑动变阻器的阻值和对应的稳定速度。
\end{enumerate}
\twopicture[numA=甲, numB=乙, wA=8cm, wB=5.2cm, align=right]
{figs/15a.png}
{figs/15b.png}

%% answer: （1）$2\UN$；（2）$\frac{1}{3}\Ums$，见详解；（3）$R=\frac{4\sqrt{2}}{7}\UO$，$v_{1}=\sqrt{2}\Ums$，$v_{2}=\frac{\sqrt{2}}{6}\Ums$
\jdanswer{
\begin{enumerate}
	\item
	$2 \UN$

	\item
	$\frac{1}{3} \Ums$；当 $v_{0}<1 \Ums$ 时，$W=\frac{1}{9} \UJ-v_{0}^{2}$；当 $v_{0}=1 \Ums$ 时，$W=0$；当 $v_{0}>1 \Ums$ 时，$W=\frac{49}{16} \UJ-v_{0}^{2}$

	\item
	$R=\frac{4\sqrt{2}}{7} \UO$，$v_{1}=\sqrt{2} \Ums$，$v_{2}=\frac{\sqrt{2}}{6} \Ums$
\end{enumerate}
}
%% memo:
\memoanswer{
（1）滑动变阻器阻值 $R=1 \UO$，金属棒速度为 $v_{1}=0.5 \Ums$ 时，金属棒产生的感应电动势 $E_{1}=BLv_{1}$

感应电流 $I_{1}=\frac{E_{1}}{R}$

金属棒受到的安培力 $F_{1}=BI_{1}L$

解得 $F_{1}=\frac{B^{2}L^{2}v_{1}}{R}=2 \UN$

（2）调节滑动变阻器阻值 $R=1 \UO$，初速度 $v_{0}=\frac{5}{6} \Ums$，设金属棒的稳定速度为 $v$，同理可得安培力 $F_{\text{安}}=\frac{B^{2}L^{2}v}{R}$

金属棒的初动能 $E_{k0}=\frac{1}{2}mv_{0}^{2}=\frac{25}{36} \UJ<2 \UJ$

根据图乙可知外力 $F=F_{0}+kE_{k}$

其中 $F_{0}=1 \UN$，$k=\frac{7-1}{2-0} \UN/\UJ=3 \UN/\UJ$

金属棒稳定时可得 $F=F_{\text{安}}$

解得 $v=\frac{1}{3} \Ums$，或 $v=1 \Ums$

初始时，安培力 $F_{\text{安}0}=\frac{B^{2}L^{2}v_{0}}{R}=\frac{10}{3} \UN$

外力 $F_{1}=F_{0}+kE_{k0}=\frac{37}{12} \UN$

$F_{1}<F_{\text{安}0}$，金属棒减速运动，故 $v=1 \Ums$ 舍去，金属棒的稳定速度 $v=\frac{1}{3} \Ums$

根据动能定理得合外力做功 $W=E_{k}-E_{k0}$

若 $v_{0}=\frac{1}{3} \Ums$ 或 $v_{0}=1 \Ums$，金属棒受到的外力等于金属棒受到的安培力，即 $F=\frac{B^{2}L^{2}v_{0}}{R}$

金属棒直接匀速直线运动，合外力做功 $W=0$

若 $v_{0}<\frac{1}{3} \Ums$，金属棒受到的外力大于金属棒受到的安培力，即 $F>\frac{B^{2}L^{2}v_{0}}{R}$

金属棒逐渐加速到稳定速度 $v=\frac{1}{3} \Ums$，合外力做功 $W=\frac{1}{9} \UJ-\frac{1}{2}mv_{0}^{2}=\frac{1}{9} \UJ-v_{0}^{2}$

若 $\frac{1}{3} \Ums<v_{0}<1 \Ums$，金属棒受到的安培力大于外力，即 $F_{0}<\frac{B^{2}L^{2}v_{0}}{R}$

金属棒逐渐减速到稳定速度 $v=\frac{1}{3} \Ums$，合外力做功 $W=\frac{1}{9} \UJ-\frac{1}{2}mv_{0}^{2}=\frac{1}{9} \UJ-v_{0}^{2}$

若 $\sqrt{2} \Ums>v_{0}>1 \Ums$，金属棒受到的外力大于金属棒受到的安培力，即 $F>\frac{B^{2}L^{2}v_{0}}{R}$

金属棒逐渐加速到 $v_{2}=\sqrt{2} \Ums$，外力恒定 $F'=7 \UN$

安培力 $F_{\text{安}}'=\frac{B^{2}L^{2}v_{2}}{R}=4\sqrt{2} \UN$，外力仍大于安培力，导体棒继续加速，设稳定时速度为 $v'$

根据 $F'=\frac{B^{2}L^{2}v'}{R}=7 \UN$

可得稳定速度 $v'=\frac{7}{4} \Ums$

可知若 $\frac{7}{4} \Ums>v_{0}>1 \Ums$，金属棒受到的外力大于金属棒受到的安培力，金属棒加速到稳定速度 $v'=\frac{7}{4} \Ums$，合外力做功 $W=\frac{49}{16} \UJ-\frac{1}{2}mv_{0}^{2}=\frac{49}{16} \UJ-v_{0}^{2}$

若 $v_{0}>\frac{7}{4} \Ums$，金属棒受到的外力小于金属棒受到的安培力，金属棒减速到稳定速度 $v'=\frac{7}{4} \Ums$，合外力做功 $W=\frac{49}{16} \UJ-v_{0}^{2}$

综上可得若 $v_{0}<1 \Ums$，合外力做功 $W=\frac{1}{9} \UJ-v_{0}^{2}$

若 $v_{0}=1 \Ums$，合外力做功 $W=0$

若 $v_{0}>1 \Ums$，合外力做功 $W=\frac{49}{16} \UJ-v_{0}^{2}$

（3）调节滑动变阻器阻值，使得金属棒存在两个稳定速度，根据稳定条件 $F=F_{\text{安}}$

当 $0\leq v_{0}\leq\sqrt{2} \Ums$ 时，可得方程 $\frac{B^{2}L^{2}v}{R}=F_{0}+k\cdot E_{k}$

其中 $E_{k}=\frac{1}{2}mv^{2}$

整理得 $\frac{km}{2}v^{2}-\frac{B^{2}L^{2}}{R}\cdot v+F_{0}=0$

根据 $\Delta=\left(\frac{B^{2}L^{2}}{R}\right)^{2}-\frac{4kmF_{0}}{2}>0$

可得 $R<\frac{2\sqrt{3}}{3} \UO$

此时两个根满足 $v_{1}+v_{2}=\frac{B^{2}L^{2}}{2kR}>0$，$v_{1}v_{2}=\frac{F_{0}}{k}>0$

可知两根都为正根，要恰好存在两个稳定速度，则其中一个较大根的端点要恰好位于 $v_{1}=\sqrt{2} \Ums$

解得此时 $R=\frac{4\sqrt{2}}{7} \UO$

解得另一个较小根 $v_{2}=\frac{\sqrt{2}}{6} \Ums$

当 $v\geq\sqrt{2} \Ums$ 时，可得方程 $\frac{B^{2}L^{2}v}{R}=7 \UN$

解得 $R\geq\frac{4\sqrt{2}}{7} \UO$

当电阻 $R=\frac{4\sqrt{2}}{7} \UO$ 时，根据外力恒定段方程 $\frac{B^{2}L^{2}v}{R}=7 \UN$

可得 $v=\sqrt{2} \Ums$

两段解重合，这里存在一个稳定速度 $v_{1}=\sqrt{2} \Ums$。

故滑动变阻器的阻值 $R=\frac{4\sqrt{2}}{7} \UO$ 时，金属棒存在两个对应的稳定速度为 $v_{1}=\sqrt{2} \Ums$ 和 $v_{2}=\frac{\sqrt{2}}{6} \Ums$。
}
\end{enumerate}

\end{document}
